\section{Класове непрекъснати случайни величини}

В тази част от лекцията ще разгледаме три основни класа непрекъснати случайни величини (равномерно, нормално и експоненциално разпределение), които имат важна роля в теорията на вероятностите, и ще покажем как структурно могат да бъдат свързани със своите стандартизирани версии чрез подходящи трансформации. Тази техника значително улеснява пресмятанията на техните моменти.

\subsection{Равномерно разпределение}

Казваме, че случайната величина $X$ има равномерно разпределение в интервала $[a, b]$, и записваме $X \sim U(a, b)$, където $-\infty < a < b < \infty$, ако нейната плътност се задава с:
\[
f_X(x) = 
\begin{cases} 
\frac{1}{b-a}, & \text{за } x \in [a, b] \\ 
0, & \text{иначе} 
\end{cases}
\]

Една от важните роли на равномерното разпределение е, че когато знаем, че възможните стойности на даден експеримент лежат в интервала $[a, b]$, но нямаме допълнителна информация, то избирайки това разпределение ние внасяме най-малко субективизъм (принцип на максималната ентропия). Всеки подинтервал с една и съща дължина има една и съща вероятност.

Функцията на разпределение $F_X(x)$ е:
\[
F_X(x) = 
\begin{cases} 
0, & \text{за } x \le a \\ 
\frac{x-a}{b-a}, & \text{за } a < x < b \\ 
1, & \text{за } x \ge b 
\end{cases}
\]

Очакваната стойност се пресмята стандартно като център на тежестта:
\[
\mathbb{E}[X] = \int_{a}^{b} x \frac{1}{b-a} dx = \frac{a+b}{2}
\]

\textbf{Структурна връзка със стандартно равномерно разпределение:} 
Вместо да пресмятаме сложни интеграли за всяко конкретно $a$ и $b$, можем да сведем $X$ до стандартното равномерно разпределение $Y \sim U(0, 1)$. Ако разгледаме трансформацията $y = g(x) = \frac{x-a}{b-a}$, която е строго монотонно растяща, то обратната ѝ функция е $x = g^{-1}(y) = a + (b-a)y$. 

Ако дефинираме $Z = g(X) = \frac{X-a}{b-a}$, по теоремата за смяна на променливите плътността на $Z$ се дава от:
\[
f_Z(y) = f_X(g^{-1}(y)) \left| \frac{d}{dy} g^{-1}(y) \right| = f_X(a + (b-a)y) \cdot (b-a)
\]
Понеже $y \in [0, 1] \implies a + (b-a)y \in [a, b]$, в този интервал $f_X$ е $\frac{1}{b-a}$. Следователно:
\[
f_Z(y) = \frac{1}{b-a} \cdot (b-a) = 1 \quad \text{за } y \in [0, 1]
\]
Това доказва, че $Z \sim U(0, 1)$. 

Чрез обратната връзка $X = a + (b-a)Z$ можем лесно да изчислим моментите на $X$. Очакването е:
\[
\mathbb{E}[X] = \mathbb{E}[a + (b-a)Z] = a + (b-a)\mathbb{E}[Z] = a + (b-a)\frac{1}{2} = \frac{a+b}{2}
\]
Дисперсията е:
\[
\text{Var}(X) = \text{Var}(a + (b-a)Z) = (b-a)^2 \text{Var}(Z)
\]
За $Z \sim U(0,1)$, вторият момент е $\mathbb{E}[Z^2] = \int_0^1 z^2 \cdot 1 dz = \frac{1}{3}$. Тогава $\text{Var}(Z) = \frac{1}{3} - \left(\frac{1}{2}\right)^2 = \frac{1}{12}$. Оттук получаваме:
\[
\text{Var}(X) = \frac{(b-a)^2}{12}
\]

\subsection{Нормално разпределение (Гаусово)}

Случайната величина $X$ има нормално разпределение с параметри $\mu$ (средно) и $\sigma^2$ (дисперсия, $\sigma > 0$), което се записва като $X \sim N(\mu, \sigma^2)$, ако нейната плътност е:
\[
f_X(x) = \frac{1}{\sqrt{2\pi}\sigma} e^{-\frac{(x-\mu)^2}{2\sigma^2}}, \quad \forall x \in \mathbb{R}
\]
Тази крива, известна като камбанка (bell-shaped curve), има огромно значение в природата и се появява естествено благодарение на Централната гранична теорема (ЦГТ). Интегралът от тази плътност над цялата реална ос е равен на 1.

\begin{defn}[стандартно нормално разпределение]\label{def:std-normal}
Специалният случай, когато $\mu = 0$ и $\sigma^2 = 1$, се нарича \emph{стандартно нормално разпределение} и се бележи със $Z \sim N(0, 1)$. Неговата плътност е
\[
f_Z(x) = \frac{1}{\sqrt{2\pi}} e^{-\frac{x^2}{2}} .
\]
\end{defn}
Функцията на разпределение на стандартното нормално разпределение няма затворена форма чрез елементарни функции и традиционно се означава с главната буква $\Phi(x)$:
\[
\Phi(x) = \int_{-\infty}^{x} \frac{1}{\sqrt{2\pi}} e^{-\frac{y^2}{2}} dy
\]

\textbf{Стандартизиране на нормална случайна величина:} 
Нека $X \sim N(\mu, \sigma^2)$. Въвеждаме трансформацията $Y = g(X) = \frac{X-\mu}{\sigma}$. Обратната функция е $x = g^{-1}(y) = \sigma y + \mu$ с производна $\sigma$.
Плътността на $Y$ е:
\[
f_Y(y) = f_X(\sigma y + \mu) \cdot \sigma = \frac{1}{\sqrt{2\pi}\sigma} e^{-\frac{(\sigma y + \mu - \mu)^2}{2\sigma^2}} \cdot \sigma = \frac{1}{\sqrt{2\pi}} e^{-\frac{y^2}{2}}
\]
Това съвпада с плътността на $Z$, следователно $\frac{X-\mu}{\sigma} \sim N(0, 1)$. Обратната връзка е $X = \sigma Z + \mu$.

Тази структурна връзка позволява лесно изчисляване на моментите. Очакването е:
\[
\mathbb{E}[X] = \sigma\mathbb{E}[Z] + \mu
\]
Тъй като подинтегралната функция $y e^{-y^2/2}$ за очакването на $Z$ е нечетна и се интегрира в симетрични граници $(-\infty, \infty)$, следва че $\mathbb{E}[Z] = 0$. Следователно $\mathbb{E}[X] = \mu$.

За дисперсията имаме $\text{Var}(X) = \sigma^2 \text{Var}(Z) = \sigma^2 \mathbb{E}[Z^2]$. За да пресметнем втория момент $\mathbb{E}[Z^2]$, ще използваме елегантен метод на Ричард Файнман — диференциране под знака на интеграла по параметър. 
Знаем, че интегралът от общата нормална плътност е 1, т.е.:
\[
\sqrt{2\pi}\sigma = \int_{-\infty}^{\infty} e^{-\frac{y^2}{2\sigma^2}} dy
\]
Диференцираме двете страни по параметъра $\sigma$:
\[
\sqrt{2\pi} = \int_{-\infty}^{\infty} \frac{\partial}{\partial \sigma} \left( e^{-\frac{y^2}{2\sigma^2}} \right) dy = \int_{-\infty}^{\infty} e^{-\frac{y^2}{2\sigma^2}} \left(-\frac{y^2}{2}\right) \left(-\frac{2}{\sigma^3}\right) dy
\]
Оттук получаваме равенството:
\[
\sqrt{2\pi}\sigma^3 = \int_{-\infty}^{\infty} y^2 e^{-\frac{y^2}{2\sigma^2}} dy
\]
Като заместим $\sigma = 1$, намираме втория момент на стандартното нормално разпределение:
\[
\sqrt{2\pi} = \int_{-\infty}^{\infty} y^2 e^{-\frac{y^2}{2}} dy \implies \mathbb{E}[Z^2] = \frac{1}{\sqrt{2\pi}} \int_{-\infty}^{\infty} y^2 e^{-\frac{y^2}{2}} dy = \frac{\sqrt{2\pi}}{\sqrt{2\pi}} = 1
\]
Следователно $\text{Var}(Z) = 1$, а дисперсията на първоначалната величина е $\text{Var}(X) = \sigma^2$.

За пресмятане на вероятности с общо нормално разпределение, свеждаме до $\Phi(x)$:
\[
\mathbb{P}(X < x) = \mathbb{P}\left(\frac{X-\mu}{\sigma} < \frac{x-\mu}{\sigma}\right) = \Phi\left(\frac{x-\mu}{\sigma}\right)
\]
Например, $\Phi(0) = \frac{1}{2}$ поради симетрията. Също така класическа стойност е $\Phi(1{,}96) \approx 0{,}975$, което означава, че $95\%$ от вероятностната маса е в интервала $[-1{,}96, 1{,}96]$.

\subsection{Експоненциално разпределение}

Казваме, че $X$ е експоненциално разпределена с параметър $\lambda > 0$, записвано $X \sim \text{Exp}(\lambda)$, ако нейната плътност е:
\[
f_X(x) = 
\begin{cases} 
\lambda e^{-\lambda x}, & \text{за } x \ge 0 \\ 
0, & \text{иначе} 
\end{cases}
\]

Функцията на разпределение е:
\[
F_X(x) = \int_{0}^{x} \lambda e^{-\lambda y} dy = 1 - e^{-\lambda x}
\]
Често е по-удобно да се работи с опашката на разпределението (tail probability), която се бележи с $\bar{F}_X(x)$:
\[
\bar{F}_X(x) = \mathbb{P}(X \ge x) = 1 - F_X(x) = e^{-\lambda x}
\]

Този клас разпределения е фундаментален заради \textbf{свойството си на безпаметност (memorylessness)}. То отразява феномена, при който ако един процес/компонент е работил $x$ време, вероятността да продължи да работи още $y$ време е същата, сякаш е започнал да работи отначало (чисто нов). Формално:
\begin{align*}
\mathbb{P}(X \ge x+y \mid X \ge x)
  &= \frac{\mathbb{P}(X \ge x+y \text{ и } X \ge x)}{\mathbb{P}(X \ge x)}
   = \frac{\mathbb{P}(X \ge x+y)}{\mathbb{P}(X \ge x)} \\
  &= \frac{e^{-\lambda(x+y)}}{e^{-\lambda x}} = e^{-\lambda y} = \mathbb{P}(X \ge y)
\end{align*}

Обратно, ако търсим функция $\bar{F}$, която да удовлетворява функционалното уравнение за безпаметност $\bar{F}(x+y) = \bar{F}(x)\bar{F}(y)$, и изискваме отговарящото ѝ разпределение да притежава плътност, единственото възможно решение е експоненциалната функция $\bar{F}(x) = e^{-\lambda x}$.

Структурната връзка тук се изразява чрез факта, че ако $X \sim \text{Exp}(\lambda)$, то мащабираната величина $Y = \lambda X \sim \text{Exp}(1)$. Проверката е директна:
\[
\mathbb{P}(Y \le x) = \mathbb{P}(\lambda X \le x) = \mathbb{P}\left(X \le \frac{x}{\lambda}\right) = 1 - e^{-\lambda \left(\frac{x}{\lambda}\right)} = 1 - e^{-x} = F_Y(x)
\]
Тъй като $Y$ е с параметър $1$, нейните моменти са $\mathbb{E}[Y] = 1$ и $\text{Var}(Y) = 1$. Чрез линейните свойства на очакването намираме за $X$:
\[
\mathbb{E}[X] = \mathbb{E}\left[\frac{Y}{\lambda}\right] = \frac{1}{\lambda}\mathbb{E}[Y] = \frac{1}{\lambda}
\]
\[
\text{Var}(X) = \text{Var}\left(\frac{Y}{\lambda}\right) = \frac{1}{\lambda^2}\text{Var}(Y) = \frac{1}{\lambda^2}
\]

\section{Непрекъснати съвместни разпределения}

В практиката най-често се работи със съвкупност от случайни величини, които може да имат сложни зависимости помежду си. Ето защо въвеждаме случайни вектори.

\begin{defn}
Векторът $X = (X_1, \dots, X_n)$ е непрекъснат вектор от случайни величини, ако съществува функция $f_X: \mathbb{R}^n \to [0, \infty)$, наречена \textbf{съвместна плътност} (joint probability density function), такава че:
1. $f_X(x) \ge 0 \quad \forall x \in \mathbb{R}^n$
2. $\int_{-\infty}^{\infty} \dots \int_{-\infty}^{\infty} f_X(x_1, \dots, x_n) dx_1 \dots dx_n = 1$
3. За всяко подмножество $D \subseteq \mathbb{R}^n$: 
   \[
   \mathbb{P}(X \in D) = \int \dots \int_D f_X(x) dx
   \]
\end{defn}

Множеството, върху което съвместната плътност е строго положителна, се нарича дефиниционно множество:
\[
\mathcal{D}(f_X) = \{ x \in \mathbb{R}^n : f_X(x) > 0 \}
\]

Ако знаем съвместната плътност на целия вектор, можем да възстановим плътността на всеки индивидуален компонент $X_j$.

\begin{defn}[маргинална плътност]\label{def:marginal}
\emph{Маргинална плътност} на компонента $X_j$ наричаме
\[
f_{X_j}(x_j) = \int_{-\infty}^{\infty} \dots \int_{-\infty}^{\infty} f_X(x_1, \dots, x_n) \, dx_1 \dots dx_{j-1} dx_{j+1} \dots dx_n ,
\]
тоест резултатът от „дезинтегриране“ по всички останали променливи.
\end{defn}

\textbf{Съвместна функция на разпределение:}
Тя се дефинира общо (не само за непрекъснати величини) като:
\[
F_X(x) = \mathbb{P}(X_1 \le x_1, \dots, X_n \le x_n) = \int_{-\infty}^{x_1} \dots \int_{-\infty}^{x_n} f_X(y_1, \dots, y_n) dy_1 \dots dy_n
\]
От функцията на разпределение може да се възстанови плътността чрез $n$-кратна смесена производна:
\[
f_X(x_1, \dots, x_n) = \frac{\partial^n}{\partial x_1 \dots \partial x_n} F_X(x_1, \dots, x_n)
\]

\textbf{Независимост:}
Две случайни величини $X_1$ и $X_2$ са независими ($X_1 \perp \!\!\! \perp X_2$), тогава и само тогава, когато съвместната им функция на разпределение се разпада на произведение от индивидуалните:
\[
F_X(x_1, x_2) = F_{X_1}(x_1) F_{X_2}(x_2)
\]
В случая на непрекъснат вектор, това е еквивалентно на разпадане на съвместната плътност:
\[
f_X(x_1, x_2) = f_{X_1}(x_1) f_{X_2}(x_2)
\]

Случайните величини $X_1, \dots, X_n$ са \textbf{независими в съвкупност}, ако за всяко $k$ ($2 \le k \le n$) и за всеки избор на $k$ различни индекса $j_1, \dots, j_k$, съвместната плътност на съответния подвектор е произведение от маргиналните плътности:
\[
f_Y(y_{j_1}, \dots, y_{j_k}) = \prod_{i=1}^{k} f_{X_{j_i}}(y_{j_i})
\]
В частност, плътността на целия вектор е $f_X(x_1, \dots, x_n) = \prod_{i=1}^n f_{X_i}(x_i)$.

\textbf{Свойства на очакването:}
Ако $X$ е непрекъснат вектор и $g$ е функция, то математическото очакване на $g(X)$ се задава като:
\[
\mathbb{E}[g(X)] = \int_{-\infty}^{\infty} \dots \int_{-\infty}^{\infty} g(x) f_X(x) dx
\]
Очакването е винаги линеен оператор. Ако $X = (X_1, X_2)$, то за функцията $g(X_1, X_2) = X_1 + X_2$:
\[
\mathbb{E}[X_1 + X_2] = \int_{-\infty}^{\infty} \int_{-\infty}^{\infty} (x_1 + x_2) f_X(x_1, x_2) dx_1 dx_2
\]
Можем да разделим интеграла на две части. Като интегрираме първата част по $x_2$, получаваме маргиналната плътност на $X_1$:
\[
\int_{-\infty}^{\infty} x_1 \left( \int_{-\infty}^{\infty} f_X(x_1, x_2) dx_2 \right) dx_1 = \int_{-\infty}^{\infty} x_1 f_{X_1}(x_1) dx_1 = \mathbb{E}[X_1]
\]
Аналогично, втората част е $\mathbb{E}[X_2]$. Следователно $\mathbb{E}[X_1 + X_2] = \mathbb{E}[X_1] + \mathbb{E}[X_2]$. 

Ако допълнително $X_1$ и $X_2$ са независими, то тогава:
\[
\mathbb{E}[X_1 X_2] = \mathbb{E}[X_1] \mathbb{E}[X_2]
\]
\[
\text{Var}(X_1 + X_2) = \text{Var}(X_1) + \text{Var}(X_2)
\]

\section{Смяна на променливите при случайни вектори}

Често се налага да намерим разпределението на нов случаен вектор $Y$, който е функция от познат вектор $X$.

\begin{thm}[Смяна на променливите]
Нека $X = (X_1, X_2)$ е непрекъснат вектор със съвместна плътност $f_X$ и дефиниционна област $\mathcal{D}(f_X)$. 
Нека $g: \mathcal{D}(f_X) \to \mathbb{R}^2$ е взаимно еднозначно изображение, т.е. $Y = g(X) = (g_1(X), g_2(X))$.
Обратната функция бележим с $X = g^{-1}(Y) = h(Y) = (h_1(Y), h_2(Y))$. 
Нека $g$ и $h$ са непрекъснати и диференцируеми в съответните си области $\mathcal{D}(f_X)$ и $g(\mathcal{D}(f_X))$.
Тогава съвместната плътност на $Y$ съществува и е равна на:
\[
f_Y(y) = f_X(h(y)) |J(y)|, \quad \forall y \in g(\mathcal{D}(f_X))
\]
където $J(y)$ е якобианата (детерминантата на матрицата от производни):
\[
J(y) = \det \begin{pmatrix} 
\frac{\partial h_1}{\partial y_1} & \frac{\partial h_1}{\partial y_2} \\ 
\frac{\partial h_2}{\partial y_1} & \frac{\partial h_2}{\partial y_2} 
\end{pmatrix}
\]
\end{thm}

Доказателството стъпва директно върху теоремата за смяна на променливите при двукратни интеграли от математическия анализ, като се запише вероятността $\mathbb{P}(Y \in A) = \mathbb{P}(X \in h(A))$.

\textbf{Обща схема за сума на независими случайни величини:}
Нека $X_1$ и $X_2$ са независими, откъдето $f_X(x_1, x_2) = f_{X_1}(x_1) f_{X_2}(x_2)$. Интересуваме се от плътността на сумата $Y_2 = X_1 + X_2$.
За да приложим теоремата, допълваме с тривиална променлива $Y_1 = X_1$:
\[
Y = \begin{pmatrix} Y_1 \\ Y_2 \end{pmatrix} = \begin{pmatrix} X_1 \\ X_1 + X_2 \end{pmatrix} = g(X)
\]
Обратната трансформация е лесна: $X_1 = Y_1$ и $X_2 = Y_2 - Y_1$. Следователно:
\[
X = h(Y) = \begin{pmatrix} Y_1 \\ Y_2 - Y_1 \end{pmatrix}
\]
Якобианата е $J(y) = \det \begin{pmatrix} 1 & 0 \\ -1 & 1 \end{pmatrix} = 1$.
Съвместната плътност на $Y$ става:
\[
f_Y(y_1, y_2) = f_X(h(y_1, y_2)) \cdot 1 = f_{X_1}(y_1) f_{X_2}(y_2 - y_1)
\]
За да намерим търсената маргинална плътност на сумата $Y_2$, интегрираме по $y_1$:

\begin{defn}[конволюция]\label{def:convolution}
Плътността на сумата на две независими непрекъснати случайни величини се дава от \emph{конволюцията}
\[
f_{Y_2}(y_2) = \int f_{X_1}(y_1) f_{X_2}(y_2 - y_1) \, dy_1 .
\]
\end{defn}

Тънкостта при прилагането ѝ е правилното определяне на границите на интеграла, тъй като те зависят от дефиниционните области на изходните величини (виж следващия пример).

\section{Гама разпределение и Хи-квадрат разпределение}

\subsection{Гама разпределение}

Казваме, че $X$ има Гама разпределение с параметри $\alpha > 0$ и $\beta > 0$, записвано $X \sim \Gamma(\alpha, \beta)$, ако плътността ѝ е:
\[
f_X(x) = 
\begin{cases}
\frac{\beta^\alpha}{\Gamma(\alpha)} x^{\alpha - 1} e^{-\beta x}, & \text{за } x \ge 0 \\
0, & \text{иначе}
\end{cases}
\]
където $\Gamma(\alpha) = \int_0^\infty x^{\alpha-1} e^{-x} dx$ е добре познатата Гама функция. Важни нейни стойности са $\Gamma(1) = 1$ и $\Gamma(n+1) = n!$.

Забележете връзката с експоненциалното разпределение: ако фиксираме $\alpha = 1$, плътността става $f_X(x) = \frac{\beta^1}{1} x^{0} e^{-\beta x} = \beta e^{-\beta x}$. Това означава, че $\Gamma(1, \beta) \equiv \text{Exp}(\beta)$.

\begin{prop}[Сума на Гама величини]
Нека $X_1, \dots, X_n$ са независими в съвкупност случайни величини, като $X_i \sim \Gamma(\alpha_i, \beta)$ (имат общ мащабен параметър $\beta$). Тогава тяхната сума също има Гама разпределение:
\[
Y = \sum_{j=1}^n X_j \sim \Gamma\left(\sum_{i=1}^n \alpha_i, \beta\right)
\]
\end{prop}

\begin{proof}[Доказателство (за $n=2$)]
Използваме конволюционната схема от предишния раздел. Търсим плътността на $Y_2 = X_1 + X_2$. Тъй като $X_1 \ge 0$ и $X_2 \ge 0$, то $y_1 \ge 0$ и $y_2 - y_1 \ge 0 \implies 0 \le y_1 \le y_2$. Границите на интеграла са от $0$ до $y_2$:
\[
f_{Y_2}(y_2) = \int_{0}^{y_2} f_{X_1}(y_1) f_{X_2}(y_2 - y_1) dy_1
\]
Заместваме плътностите на двете Гама разпределения:
\[
f_{Y_2}(y_2) = \int_{0}^{y_2} \left( \frac{\beta^{\alpha_1}}{\Gamma(\alpha_1)} y_1^{\alpha_1 - 1} e^{-\beta y_1} \right) \left( \frac{\beta^{\alpha_2}}{\Gamma(\alpha_2)} (y_2 - y_1)^{\alpha_2 - 1} e^{-\beta (y_2 - y_1)} \right) dy_1
\]
Експоненциалните членове се опростяват: $e^{-\beta y_1} e^{-\beta y_2 + \beta y_1} = e^{-\beta y_2}$, което не зависи от $y_1$ и излиза пред интеграла. Остава да се реши:
\[
f_{Y_2}(y_2) = \frac{\beta^{\alpha_1+\alpha_2}}{\Gamma(\alpha_1)\Gamma(\alpha_2)} e^{-\beta y_2} \int_{0}^{y_2} y_1^{\alpha_1 - 1} (y_2 - y_1)^{\alpha_2 - 1} dy_1
\]
С помощта на субституцията $y_1 = u y_2$, интегралът се свежда до Бета функцията $B(\alpha_1, \alpha_2) = \frac{\Gamma(\alpha_1)\Gamma(\alpha_2)}{\Gamma(\alpha_1+\alpha_2)}$ и крайният резултат е точно плътността на $\Gamma(\alpha_1+\alpha_2, \beta)$:
\[
f_{Y_2}(y_2) = \frac{\beta^{\alpha_1+\alpha_2}}{\Gamma(\alpha_1+\alpha_2)} y_2^{\alpha_1+\alpha_2-1} e^{-\beta y_2}
\]
От това следва, че Гама разпределението е \textbf{безгранично делимо}. Например $\Gamma(3{,}5, \beta)$ може да се представи като сума от 7 независими $\Gamma(0{,}5, \beta)$ величини.
\end{proof}

\subsection{Хи-квадрат разпределение ($\chi^2$)}

\begin{prop}
Нека $Z_1, \dots, Z_n$ са независими в съвкупност стандартно нормални случайни величини ($Z_i \sim N(0, 1)$). Тогава сумата от техните квадрати има $\chi^2$ разпределение с $n$ степени на свобода, което е еквивалентно на Гама разпределение:
\[
Y = \sum_{j=1}^n Z_j^2 \sim \chi^2(n) \equiv \Gamma\left(\frac{n}{2}, \frac{1}{2}\right)
\]
\end{prop}

\begin{proof}
Ще покажем, че квадратът на една стандартно нормална величина $T_1 = Z_1^2$ има разпределение $\chi^2(1) \equiv \Gamma(1/2, 1/2)$. 
Започваме с функцията на разпределение за $t > 0$:
\[
F_{T_1}(t) = \mathbb{P}(T_1 < t) = \mathbb{P}(Z_1^2 < t) = \mathbb{P}(-\sqrt{t} < Z_1 < \sqrt{t}) = \frac{1}{\sqrt{2\pi}} \int_{-\sqrt{t}}^{\sqrt{t}} e^{-\frac{y^2}{2}} dy
\]
Поради четността на функцията, това е:
\[
F_{T_1}(t) = 2 \cdot \frac{1}{\sqrt{2\pi}} \int_{0}^{\sqrt{t}} e^{-\frac{y^2}{2}} dy
\]
Диференцираме по $t$, за да намерим плътността (използвайки правилото за диференциране на интеграл с променлива горна граница $\sqrt{t}$, чиято производна е $\frac{1}{2\sqrt{t}}$):
\[
f_{T_1}(t) = \frac{d}{dt} F_{T_1}(t) = \frac{2}{\sqrt{2\pi}} e^{-\frac{(\sqrt{t})^2}{2}} \cdot \frac{1}{2\sqrt{t}} = \frac{1}{\sqrt{2\pi}} t^{-1/2} e^{-t/2}
\]

Сега ще сравним този резултат с плътността на Гама разпределението при параметри $\alpha = 1/2$ и $\beta = 1/2$:
\[
f_\Gamma(t) = \frac{(1/2)^{1/2}}{\Gamma(1/2)} t^{(1/2) - 1} e^{-t/2} = \frac{1}{\sqrt{2}\Gamma(1/2)} t^{-1/2} e^{-t/2}
\]
Тъй като и двата израза представляват плътност (и следователно интегралите им са 1), константите отпред трябва да са равни. От равенството $\frac{1}{\sqrt{2\pi}} = \frac{1}{\sqrt{2}\Gamma(1/2)}$ можем пътем да изчислим красивия факт, че $\Gamma(1/2) = \sqrt{\pi}$.
Най-важното е, че двете плътности съвпадат. Следователно $Z_1^2 \sim \Gamma(1/2, 1/2)$.

Тъй като $Y = \sum_{j=1}^n Z_j^2$ е сума на $n$ независими Гама величини с еднакъв мащабен параметър $\beta = 1/2$ и формов параметър $\alpha_j = 1/2$, от предходното твърдение следва директно, че:
\[
Y \sim \Gamma\left(\sum_{j=1}^n \frac{1}{2}, \frac{1}{2}\right) = \Gamma\left(\frac{n}{2}, \frac{1}{2}\right) \equiv \chi^2(n)
\]
С това доказателството е завършено.
\end{proof}
